\BourbakiStarSection{Historical Note}

(Numbers in brackets refer to the bibliography at the end of this Note.)

There are few notions in Mathematics more primitive than that of a law of composition: it seems inseparable from the first rudiments of calculations on natural numbers and measurable quantities. The most ancient documents that have come down to us on the Mathematics of the Egyptians and Babylonians show that they already had a complete system of rules for calculating with natural numbers \textgreater0, rational numbers \textgreater0, lengths and areas; although the preserved texts deal only with problems in which the given quantities have explicit numerical values,† they leave us in no doubt concerning the generality they attributed to the rules employed and show a quite remarkable technical ability in the manipulation of equations of the first and second degree ({[}1{]}, p.~179 et seq.). On the other hand there is not the slightest trace of a concern to justify the rules used nor even of precise definitions of the operations which occur: the latter as the former arise in a purely empirical manner.

Just such a concern, however, is already very clearly manifest in the works of the Greeks of the classical age; true, an axiomatic treatment is not found of the theory of natural numbers (such an axiomatization was only to appear at the end of the nineteenth century; see the Historical Note to Set Theory, IV); but in Euclid's Elements there are numerous passages giving formal proofs of rules of calculation which are just as intuitively ``obvious'' as those of calculation with integers (for example, the commutativity of the product of two rational numbers). The most remarkable proofs of this nature are those relating to the theory of magnitudes, the most original creation of Greek Mathematics (equivalent, as is known, to our theory of real numbers \textgreater0; see the Historical Note to General Topology, IV); here Euclid considers amongst other things the product of two ratios of magnitudes and shows that it is independent of the form of presentation of these ratios (the first example of a ``quotient'' of a law of composition by an equivalence relation in the sense of § 1, no. 6) and that they commute ({[}2{]}, Book V, Prop. 22--23).†

It must not however be concealed that this progress towards rigour is accompanied in Euclid by a stagnation and even in some ways by a retrogression as far as the technique of algebraic calculations is concerned. The overwhelming preponderance of Geometry (in the light of which the theory of magnitudes was obviously conceived) paralyses any independent development of algebraic notation: the elements entering into the calculations must always be ``represented'' geometrically; moreover the two laws of composition which occur are not defined on the same set (addition of ratios is not defined in a general way and the product of two lengths is not a length but an area); the result is a lack of flexibility making it almost impracticable to manipulate algebraic relations of degree greater than the second.

Only with the decline of classical Greek Mathematics does Diophantus return to the tradition of the ``logisticians'' or professional calculators, who had continued to apply such rules as they had inherited from the Egyptians and Babylonians: no longer encumbering himself with geometric representations of the ``numbers'' he considers, he is naturally led to the development of rules of abstract algebraic calculation; for example, he gives rules which (in modern language) are equivalent to the formula \(x^{m+n} = x^{m}x^{n}\) for small values (positive or negative) of m and n ({[}3{]}, vol.~I, pp.~8--13); a little later (pp.~12--13), the ``rule of signs'' is stated, the beginnings of calculating with negative numbers‡; finally, Diophantus uses for the first time a letter to represent an unknown in an equation. In contrast, he seems little concerned to attach a general significance to the methods he applies for the solution of his problems; as for the axiomatic conception of laws of composition as begun by Euclid, this appears foreign to Diophantus' thought as to that of his immediate successors; it only reappears in Algebra at the beginning of the 19th century.

There were first needed, during the intervening centuries, on the one hand the development of a system of algebraic notation adequate for the expression of abstract laws and on the other a broadening of the notion of ``number'' sufficient to give it, through the observation of a wide enough variety of special cases, a general conception. For this purpose, the axiomatic theory of ratios of magnitudes created by the Greeks was insufficient, for it did no more than make precise the intuitive concept of a real number \textgreater0 and the operations on these numbers, which were already known to the Babylonians in a more confused form; now ``numbers'' would be considered of which the Greeks had had no concept and which to begin with had no sensible ``representation'': on the one hand, zero and negative numbers which appeared in the High Middle Ages in Hindu Mathematics; on the other, imaginary numbers, the creation of Italian algebraists of the 16th century.

With the exception of zero, introduced first as a separating sign before being considered as a number (see the Historical Note to Set Theory, III), the common character of these extensions is that (at least to begin with) they were purely ``formal''. By this it must be understood that the new ``numbers'' appeared first of all as the result of operations applied to situations where, keeping to their strict definition, they had no meaning (for example, the difference \(a - b\) of two natural numbers when \(a < b\) ): whence the names ``false'', ``fictive'', ``absurd'', ``impossible'', ``imaginary'', etc., numbers attributed to them. For the Greeks of the Classical Period, enamoured above all of clear thought, such extensions were inconceivable; they could only arise with calculators more disposed than were the Greeks to display a somewhat mystic faith in the power of their methods (``the generality of Analysis'' as the 18th century would say) and to allow themselves to be carried along by the mechanics of their calculations without investigating whether each step was well founded; a confidence moreover usually justified a posteriori, by the exact results to which the extension led, with these new mathematical entities, of rules of calculation uniquely valid, in all rigour, for the numbers which were already known. This explains how little by little these generalizations of the concept of number, which at first occurred only as intermediaries in a sequence of operations whose starting and finishing points were genuine ``numbers'', came to be considered for their own sakes (independent of any application to concrete calculations); once this step had been taken, more or less tangible interpretations were sought, new objects which thus acquired the right to exist in Mathematics.†

On this subject, the Hindus were already aware of the interpretation to be given to negative numbers in certain cases (a debt in a commercial problem, for example). In the following centuries, as the methods and results of Greek and Hindu mathematics spread to the West (via the Arabs), the manipulation of these numbers became more familiar and they took on other ``representations'' of a geometric or kinematic character. With the progressive improvement in algebraic notation, this was the only notable progress in Algebra during the close of the Middle Ages.

At the beginning of the 16th century, Algebra took on a new impulse thanks to the discovery by mathematicians of the Italian school of the solution ``by radicals'' of equations of the 3rd and then of the 4th degree (of which we shall speak in more detail in the Historical Note to Algebra, V); it is at this point that, in spite of their aversion, they felt compelled to introduce into their calculations imaginary numbers; moreover, little by little confidence is born in the calculus of these ``impossible'' numbers, as in that of negative numbers, even though no ``representation'' of them was conceived for more than two centuries.

On the other hand, algebraic notation was decisively perfected by Viète and Descartes; from the latter onwards, algebraic notation is more or less that which we use today.

From the middle of the 17th to the end of the 18th century, it seems that the vast horizons opened by the creation of Infinitesimal Calculus resulted in something of a neglect of Algebra in general and especially of mathematical reflection on laws of composition and on the nature of real and complex numbers.† Thus the composition of forces and the composition of velocities, well known in Mechanics from the end of the 17th century, had no repercussions in Algebra, even though they already contained the germ of vector calculus. In fact it was necessary to await the movement of ideas which, about 1800, led to the geometric representation of complex numbers (see the Historical Note to General Topology, VIII) to see addition of vectors used in Pure Mathematics.‡

It was round about the same time that for the first time in Algebra the notion of law of composition was extended in two different directions to elements which have only distant analogies with ``numbers'' (in the broadest sense of the word up till then). The first of these extensions is due to C. F. Gauss occasioned by his arithmetical researches on quadratic forms \(ax^{2} + bxy + cy^{2}\) with integer coefficients. Lagrange had defined on the set of forms with the same discriminant an equivalence relation† and had on the other hand proved an identity which provided this set with a commutative law of composition (not everywhere defined); starting from these results, Gauss showed that this law is compatible (in the sense of § 1, no. 6) with a refined form of the above equivalence relation ({[}5{]}, vol.~I, p.~272): ``This shows'', he then says, ``what must be understood by the composite class of two or several classes''. He then proceeds to study the ``quotient'' law, establishes essentially that it is (in modern language) a commutative group law and does so with arguments whose generality for the most part goes far beyond the special case considered by Gauss (for example, the argument by which he proves the uniqueness of the inverse element is identical with the one we have given for an arbitrary law of composition in § 2, no 3 (ibid., p.~273)). Nor does he stop there: returning a little later to this question, he recognizes the analogy between composition of classes and multiplication of integers modulo a prime number‡ (ibid., p.~371), but proves also that the group of classes of quadratic forms of given discriminant is not always a cyclic group; the indications he gives on this subject show clearly that he had recognized, at any rate in this particular case, the general structure of finite Abelian groups, which we shall study in Chapter VII ({[}5{]}, vol.~I, p.~374 and vol.~II, p.~266).

The other series of research of which we wish to speak also ended in the concept of a group and its definitive introduction into Mathematics: this is the ``theory of substitutions'', the development of the ideas of Lagrange, Vandermonde and Gauss on the solution of algebraic equations. We do not give here the detailed history of this subject (see the Historical Note to Algebra, V); we must refer to the definition by Ruffini and then Cauchy ({[}6{]}, (2), vol.~I, p.~64) of the ``product'' of two permutations of a finite set,§ and the first notions concerning finite permutation groups: transitivity, primitivity, identity element, permutable elements, etc. But these first results remain, as a whole, quite superficial and Évariste Galois must be considered as the true initiator of the theory: having, in his memorable work {[}8{]}, reduced the study of algebraic equations to that of permutation groups associated with them, he took the study of the latter considerably deeper, both with regard to the general properties of groups (Galois was the first to define the notion of normal subgroup and to recognize its importance) and with regard to the determination of groups with particular properties (where the results he obtained are still today considered to be amongst the most subtle in the theory). The first idea of the ``linear representation of groups'' also goes back to Galois† and this fact shows clearly that he possessed the notion of isomorphism of two group structures, independently of their ``realizations''.

However, if the ingenious methods of Gauss and Galois incontestably led them to a very broad conception of the notion of law of composition, they did not develop their ideas particularly on this point and their works had no immediate effect on the evolution of Abstract Algebra.‡ The clearest progress towards abstraction was made in a third direction: after reflecting on the nature of imaginary numbers (whose geometric representation had brought about, at the beginning of the 19th century, a considerable number of works), algebraists of the English school were the first, between 1830 and 1850, to isolate the abstract notion of law of composition and immediately they broadened the field of Algebra by applying this notion to a host of new mathematical entities: the algebra of Logic with Boole (see the Historical Note to Set Theory, IV), vectors and quaternions with Hamilton, general hypercomplex systems, matrices and non-associative laws with Cayley ({[}10{]}, vol.~I, pp.~127 and 301 and vol.~II, pp.~185 and 475). A parallel evolution was pursued independently on the Continent, notably in Vector Calculus (Möbius, Bellavitis), Linear Algebra and hypercomplex systems (Grassmann), of which we shall speak in more detail in the Historical Note to Algebra, III.§

From all this ferment of original and fertile ideas which revitalized Algebra in the first half of the 19th century the latter emerged thoroughly renewed. Previously methods and results had gravitated round the central problem of the solution of algebraic equations (or Diophantine equations in the Theory of Numbers): ``Algebra'', said Serret in the Introduction to his Course of Higher Algebra {[}12{]}, ``is, properly speaking, the Analysis of equations''. After 1850, if the treatises on Algebra still gave preeminence to the theory of equations for some time to come, new research works were no longer dominated by the concern for immediate applications to the solution of numerical equations but were oriented more and more towards what today is considered to be the essential problem of Algebra, the study of algebraic structures for their own sake.

These works fall quite neatly into three main currents which extend respectively the three movements of ideas which we have analysed above and pursue parallel courses without noticeably influencing one another until the last years of the 19th century.†

The first of these arose out of the work of the German school of the 19th century (Dirichlet, Kummer, Kronecker, Dedekind, Hilbert) in building up the theory of algebraic numbers, issuing out of the work of Gauss who was responsible for the first study of this type, that of the numbers \(a + bi\) ( \(a\) and \(b\) rational). We shall not follow the evolution of this theory here (see the Historical Note to Commutative Algebra, VII): we need only mention, for our purposes, the abstract algebraic notions which arose here. Right from the first successors of Gauss, the idea of a field (of algebraic numbers) is fundamental to all the works on this subject (as also to the works of Abel and Galois on algebraic equations); its field of application was enlarged when Dedekind and Weber {[}13{]} modelled the theory of algebraic functions of one variable on that of algebraic numbers. Dedekind {[}14{]} was also responsible for the notion of an ideal which provides a new example of a law of composition between sets of elements; he and Kronecker were responsible for the more and more important role played by commutative groups and modules in the theory of algebraic fields; we shall return to this in the Historical Notes to Algebra, III, V and VII.

We shall also refer in the Historical Notes to Algebra, III and VIII to the development of Linear Algebra and hypercomplex systems which was pursued without introducing any new algebraic notion during the end of the 19th and the beginning of the 20th century, in England (Sylvester, W. Clifford) and America (B. and C. S. Pierce, Dickson, Wedderburn) following the path traced by Hamilton and Cayley, in Germany (Weierstrass, Dedekind, Frobenius, Molien) and in France (Laguerre, E. Cartan) independently of the Anglo-Saxons and using quite different methods.

As far as group theory is concerned, this developed first of all mainly as the theory of finite permutation groups, following the publication of the works of Galois and their diffusion through the works of Serret {[}12{]} and above all the great ``Traité des Substitutions'' of C. Jordan {[}15{]}. The latter there developed, greatly improving them, the works of his predecessors on the particular properties of permutation groups (transitivity, primitivity, etc.), obtaining results most of which have not been superseded since; he made an equally profound study of a very important particular class of groups, linear groups and their subgroups; moreover, he it was who introduced the fundamental notion of a homomorphism of one group into another and also (a little later) that of quotient group and who proved a part of the theorem known as the ``Jordan-Hölder Theorem''.† Finally Jordan was the first to study infinite groups {[}16{]}, which S. Lie on the one hand and F. Klein and H. Poincaré on the other were to develop considerably in two different directions several years later.

In the meantime, what is essential in a group, that is its law of composition and not the nature of the objects which constitute the group, slowly come to be realized (see for example {[}10{]}, vol.~II, pp.~123 and 131 and {[}14{]}, vol.~III, p.~439). However, even the research works on finite abstract groups had for a long time been conceived as the study of permutation groups and only around 1880 was the independent theory of finite groups beginning to develop consciously. We cannot pursue further the history of this theory which is only touched on very superficially in this Treatise; we shall just mention two of the tools which are still today among the most used in the study of finite groups and which both go back to the 19th century: the Sylow theorems‡ on p-groups which date from 1872 {[}17{]} and the theory of characters created in the last years of the century by Frobenius {[}19{]}. We refer the reader who wishes to go deeply into the theory of finite groups and the many difficult problems it raises, to the monographs of Burnside {[}20{]}, Speiser {[}23{]}, Zassenhaus {[}24{]} and Gorenstein {[}27{]}.

Around 1880 too the systematic study of group presentations was begun; previously only presentations of particular groups had been encountered, for example the alternating group \(\mathfrak{A}_5\) in a work of Hamilton {[}9 bis{]} or the monodromy groups of linear differential equations and Riemann surfaces (Schwarz, Klein, Schläfli). W. Dyck was the first {[}18{]} to define (without giving it a name) the free group generated by a finite number of generators, but he was less interested in it for its own sake than as a ``universal'' object allowing a precise definition of a group given ``by generators and relations''. Developing an idea first set forth by Cayley and visibly influenced on the other hand by the works of Riemann's school mentioned above, Dyck described an interpretation of a group of given presentation, where each generator is represented by a product of two inversions with respect to tangent or intersecting circles (see also Burnside {[}20{]}, Chapter XVIII). A little later, after the development of the theory of automorphic functions by Poincaré and his successors, then the introduction by Poincaré of the tools of Algebraic Topology, the first studies of the fundamental group would be on an equal footing with those of group presentations (Dehn, Tietze), the two theories lending one another mutual support. Moreover it was a topologist, J. Nielsen, who in 1924 introduced the term free group in the first deep study of its properties {[}21{]}, almost immediately afterwards E. Artin (always with reference to questions in topology) introduced the notion of the free product of groups and O. Schreier defined more generally the free product with amalgamated subgroups {[}22{]} (cf.~I, § 7, no. 3, Exercise 20). Here again we cannot go into the history of the later developments in this direction, referring the reader to the work {[}26{]} for more details.

There is no further room to speak of the extraordinary success, since the end of the 19th century, of the idea of group (and that of invariant which is intimately related with it) in Analysis, Geometry, Mechanics and Theoretical Physics. An analogous invasion by this notion and the algebraic notions related to it (groups with operators, rings, ideals, modules) into parts of Algebra which until then seemed quite unrelated is the distinguishing mark of the last period of evolution which we retrace here and which ended with the synthesis of the three branches which we have followed above. This unification is above all the work of the German school of the years 1900--1930: begun by Dedekind and Hilbert in the last years of the 19th century, the work of axiomatization of Algebra was vigorously pursued by E. Steinitz, then, starting in 1920, under the impulse of E. Artin, E. Noether and the algebraists of their school (Hasse, Krull, O. Schreier, van der Waerden). The treatise by van der Waerden {[}25{]}, published in 1930, brought together these works for the first time in a single exposition, opening the way and serving as a guide to many research works in Algebra during more than twenty years.

